\documentclass[a4paper,12pt]{extarticle}
\usepackage[table,xcdraw]{xcolor}
\usepackage{pgfplots}
\usepackage[unicode, pdftex]{hyperref}
\usepackage[T2A]{fontenc}
\usepackage[utf8]{inputenc}
\usepackage[english,russian]{babel}
\usepackage{amsmath,amsfonts,amsthm, mathtools}
\usepackage{amssymb}
\usepackage{dsfont}
\usepackage{graphicx}
\graphicspath{{images/}}
\usepackage{wrapfig}
\usepackage{floatrow}
\usepackage{setspace}
\usepackage[headings]{fancyhdr}
\usepackage[margin=10pt,font={small,stretch=0,9},labelfont=bf,labelsep=period,
justification=centerlast]{caption}
\usepackage{array,tabularx,tabulary,booktabs}
\RequirePackage{multirow}
\RequirePackage{multicol}
\RequirePackage{longtable}
\usepackage{indentfirst}
\usepackage{hyperref}
\usepackage{icomma}
\newcommand*{\hm}[1]{#1\nobreak\discretionary{}
	{\hbox{$\mathsurround=0pt #1$}}{}}
\usepackage{soulutf8} 
\usepackage{etoolbox} 
\usepackage{geometry}
\geometry{top=20mm}
\geometry{bottom=20mm}
\geometry{left=20mm}
\geometry{right=20mm}
\DeclareMathOperator{\sgn}{\mathop{sgn}}
\usepackage[OT1]{fontenc}
\usepackage{amsmath}
\usepackage{amsfonts}
\usepackage{amssymb}
\usepackage{wasysym}
\usepackage{wrapfig}
\usepackage{physics}
\usepackage[normalem]{ulem}
\graphicspath{{pictures/}}
\DeclareGraphicsExtensions{.png,.jpg,.jpeg}
\usepackage{indentfirst}
\usepackage[T2A]{fontenc}
\usepackage{subfigure}
\usepackage{fancyhdr}
\usepackage{caption}

\usepackage{tabularx}
\usepackage{floatrow}
\usepackage{multicol}
\usepackage{lastpage}

\pgfplotsset{
  MyStyle1/.style={
    axis lines = middle,
    enlargelimits = true,
    grid=both,
    grid style={line width=.1pt, draw=gray!10},
    major grid style={line width=.2pt,draw=gray!50},
    minor tick num=5,
    enlargelimits=0.05,
    ticklabel style={font=\tiny,fill=white},
    xlabel style={at={(ticklabel* cs:1)},anchor=north west},
    ylabel style={at={(ticklabel* cs:1)},anchor=south east},
  }
}
\def\Slope{1.5}
\def\Intercept{-20}


% Настройка отчёркиваний и нумерации
\usepackage{fancyhdr}
\pagestyle{fancy}
\fancyhf{}

\fancyhead[L]{\nouppercase{\leftmark\hfill}}
\fancyfoot[RO, RE]{Страница \thepage \ из \pageref{LastPage}}
\renewcommand{\headrulewidth}{0.5pt}
\renewcommand{\footrulewidth}{0.5pt}

\usepackage{titlesec}
\titlelabel{\thetitle.\quad}

\pgfplotsset{compat=1.18}
\begin{document}

\thispagestyle{empty}
\begin{center}
	\textit{Федеральное государственное автономное образовательное\\ учреждение высшего образования }
	
	\vspace{0.5ex}
	
	\textbf{«Московский физико-технический институт\\ (национальный исследовательский университет)»}
\end{center}

\vspace{10ex}
\begin{figure}[!h]
  \begin{center}
    
\includegraphics[width = 0.8\textwidth]{MIPT.png}
  \end{center}
\end{figure}

\begin{center}
	
	\textbf{Лабораторная работа № 3.2.4}
	
	\vspace{1ex}
	
	по курсу общей физики
	
	на тему:
	
	\textbf{\textit{<<Свободные и вынужденные
 колебания в электрическом контуре>>}}
	
	\vspace{20ex}
	
	\begin{flushright}
		\noindent
		\textit{Работу выполнил:}\\  
		\textit{Легких Никита \\(группа Б05-301)}
	\end{flushright}
	\vfill
	Долгопрудный \\ \today
	\setcounter{page}{1}
\end{center}
\newpage

% Конец шаблона
\newpage

\textbf{Цель работы:} исследование свободных и вынужденных колебаний в
колебательном контуре. \\ \\
\textbf{В работе используются:} осциллограф АКТАКОМ ADS-6142H, генератор сигналов специальной формы АКИП-3409/4, магазин сопротивления МСР-60, магазин емкости Р5025, магазин индуктивности Р567 типа
МИСП, соединительная коробка с шунтирующей емкостью, соединительные одножильные и коаксиальные провода.


\section*{Теоретическая справка}


\begin{wrapfigure}{l}{0.35\linewidth} 
	\includegraphics[width=6cm]{3.2.4 Свободные и вынужденные колебания в электрическом контуре/картинки/RLC.png}
	\caption{Колебательный контур}
	\label{RLC}
\end{wrapfigure}

Основное уравнение колебательного контура 

\begin{equation}\label{ddot I}
\ddot{I} + 2\gamma\dot{I} + \omega_0^2I = 0
\end{equation}

Где $ \gamma = \dfrac{R}{2L} $ --- коэффициент затухания, $ \omega_0^2 = \dfrac{1}{LC} $ --- собственная частота контура. Решением этого уравнения являются затухающие колебания:

\begin{equation}\label{}
I = A e^{-\gamma t} \cos (\omega t - \theta)
\end{equation}

Здесь $ \omega = \sqrt{\omega_0^2 - \gamma^2} $. Можно записать решение \eqref{ddot I} и для напряжения:

\begin{equation}\label{}
U_C = U_0 \dfrac{\omega_0}{\omega} e^{-\gamma t}\cos (\omega t - \theta)
\end{equation}

В контуре со слабым затуханием $ (\omega \backsimeq \omega_0) $ верна \textbf{формула Томпсона} для периода: 

\begin{equation}\label{}
T = \dfrac{2\pi}{\omega_0} \backsimeq  \dfrac{2\pi}{\omega} = 2\pi\sqrt{LC}
\end{equation}

Режим работы контура, при котором $ \gamma = \omega_0 $, называется \textbf{критическим}. Его сопротивление равно 

\begin{equation}\label{}
R_{кр} = 2\sqrt{\dfrac{L}{C}}
\end{equation}

Потери затухающих колебаний принято характеризовать через \textbf{добротность} и \textbf{логарифмический декремент затухания}: 

\begin{equation}\label{Q}
Q = 2\pi \dfrac{W}{\Delta W} = \dfrac{1}{R} \sqrt{\dfrac{L}{C}} \quad - \quad \text{Добротность, потери энергии}
\end{equation}
\begin{equation}\label{theta}
\Theta = \dfrac{1}{n} \gamma T = \dfrac{1}{n} \ln \dfrac{U_k}{U_{k+n}}  \quad - \quad \text{Лог. декремент, потери амплитуды}
\end{equation}

\section*{Установка}
Схема установки для исследования колебаний приведена на рисунке 1.
Колебательный контур состоит из постоянной индуктивности $L$ с активным сопротивлением $R_L$, переменной емкости $C$ и сопротивления $R$. Картина колебаний напряжения на емкости наблюдается на экране двухканального осциллографа. Для возбуждения затухающих колебаний используется генератор сигналов специальной формы. Сигнал с генератора поступает через конденсатор $C_1$ на вход колебательного контура. Данная емкость необходима чтобы выходной импеданс генератора был
много меньше импеданса колебательного контура и не влиял на процессы, проходящие в контуре.
\begin{figure}[h]
\centering
\includegraphics[width=1\linewidth]{3.2.4 Свободные и вынужденные колебания в электрическом контуре/картинки/схема.png}
\label{fig:ust}
\end{figure}
Установка предназначена для исследования не только возбужденных, но и свободных колебаний в электрической цепи. При изучении свободно затухающих колебаний генератор специальных сигналов на вход колебательного контура подает периодические короткие импульсы, которые заряжают конденсатор $C$. За время между последовательными импульсами происходит разрядка конденсатора через резистор и катушку индуктивности. Напряжение на конденсаторе $U_C$ поступает на вход канала 1(X) электронного осциллографа. Для наблюдения фазовой картины затухающих колебаний на канал 2(Y) подается напряжение с резистора R (пунктирная линия на схеме установки), которое пропорционально току $I$ ($I \sim dU_C/dt$).

При изучении возбужденных колебаний на вход колебательного контура подается
синусоидальный сигнал. С помощью осциллографа возможно измерить зависимость
амплитуды возбужденных колебаний в зависимости от частоты внешнего сигнала,
из которого возможно определить добротность колебательного контура. Альтернативным способом расчета добротности контура является определение декремента
затухания по картине установления возбужденных колебаний. В этом случае генератор сигналов используется для подачи цугов синусоидальной формы.


\section*{Задание}
\section*{2.2 Измерение периодов свободных колебаний}
Устанавливаем $R = 0$ Ом, $L = 100$ Гн, $C = 0$ мкФ.

Измеряем период при помощи осциллографа (причем измеряем расстояние именно между пересечениями оси $Ot$) $T = 65.2$ мкс, отсюда по формуле (\ref{eq:parazit_C}) получаем величину паразитной емкости 
\begin{equation}
    T = 2\pi \sqrt{LC} \Rightarrow C_{\text{пар}} = C_0 = \frac{T^2}{4\pi^2 L} \approx 1.1\;\;\text{нФ}
    \label{eq:parazit_C}
\end{equation}

Также снимем зависимость периода затухающих колебаний от емкости контура: \\
\begin{tabular}{|c|c|c|}
\hline
$C$, нФ & $T_{\text{эксп}}$, мкс & $T_{\text{теор}}$, мкс   \\ \hline
1.1  & 65.2  & 65.9  \\ \hline
2.1  & 90.5  & 91.1  \\ \hline
3.1  & 112.0   & 110.6 \\ \hline
4.1  & 127.3 & 127.2 \\ \hline
5.1  & 142.4 & 141.9 \\ \hline
6.1  & 155.2 & 155.2 \\ \hline
7.1  & 168.0   & 167.4 \\ \hline
8.1  & 178.4 & 178.8 \\ \hline
9.1  & 188.8 & 189.5 \\ \hline
10.1 & 198.4 & 199.7 \\ \hline
\end{tabular}\label{tb:T(C)_zatuh_koleb}
\begin{minipage}{1\textwidth}
      \begin{tikzpicture}
        \begin{axis}[
          width=0.6\linewidth,
          height=,
          xlabel={$T_{\text{теор}}$, мкс},
          ylabel={$T_{\text{эксп}}$, мкс},
          MyStyle1,
          ]
          \draw[pink, domain=0:250] plot (\x, 1*\x);
          \addplot [only marks] plot [error bars/.cd,
                x dir=both, x fixed = 0.05,
                y dir=both, y explicit,
                      ] table [
                      x=T2, y=T1,
                x error minus=dT2, x error plus=dT2, y error minus=dT1, y error plus=dT1, col sep=comma] {csv/1.csv};
        \end{axis}
      \end{tikzpicture}
    \end{minipage}
\\
Расчет стобцов ведем по следуюзим формулам: $T_{\text{теор}} = 2\pi \sqrt{LC}$, $\Delta T = T \frac{\Delta C}{2C}$\\
Погрешности считаем учитывая что $\Delta C = 0.05$ нФ, $\Delta T_{\text{эксп}} = 1$ мкс; все погрешности < 1\% $ \Rightarrow$ в данном случае они пренебрежимо малы.

\section*{2.3 Критическое сопротивление и декремент затухания}
Устанавливаем $L = 100$ Гн. Хотим также $\nu_0 = 6.5$ кГц для этого рассчитаем $C^*$ по формуле (\ref{eq:C^*_2.3}) и выставим его на установке. Опорное значение сопротивления - критическое сопротивление контура $R_{cr}$, критическое оно в том смысле что режим колебаний переходит от периодического к апериодическому, рассчитываем по формуле (\ref{eq:R_cr_2.3}). 
\begin{equation}
    \nu_0 = \frac{1}{2\pi \sqrt{LC}} \Rightarrow C^* = \frac{1}{\nu_0^2 4 \pi^2 L} = 6\;\;\text{нФ}
    \label{eq:C^*_2.3}
\end{equation}
\begin{equation}
    R_{cr} =  \sqrt{\frac{L}{C^*}} = 8160\;\;\text{Ом}
    \label{eq:R_cr_2.3}
\end{equation}
Далее снимаем кривые напряжения $U$ при различных сопротивлениях $R$ и посчитаем добротность нашего контура:\\
\begin{tabular}{|c|c|c|c|c|c|c|c|c|}
\hline
$R$, Ом    & $U_i$, В     & $U_{i+n}$, В    & $n$ & $\Theta$ & $R_\Sigma$, Ом & $\frac{1}{\Theta^2}$ &  $\frac{1}{R_\Sigma^2} \cdot 10^{-6}$, Ом$^{-2}$ & $Q$ \\ \hline
410  & 9.80  & 4.68 & 2 & 0.37 & 487  & 7.32                   & 4.22                   & 8.50 \\ \hline
610  & 12.50 & 4.56 & 2 & 0.50 & 687  & 3.93                   & 2.12                   & 6.23 \\ \hline
810  & 11.32 & 3.10 & 2 & 0.65 & 887  & 2.38                   & 1.27                   & 4.85 \\ \hline
1010 & 9.90  & 4.60 & 1 & 0.77 & 1087 & 1.70                   & 0.85                   & 4.10 \\ \hline
1210 & 9.04  & 3.52 & 1 & 0.94 & 1287 & 1.12                   & 0.60                   & 3.33 \\ \hline
1410 & 8.08  & 2.72 & 1 & 1.09 & 1487 & 0.84                   & 0.45                   & 2.89 \\ \hline
1610 & 7.32  & 1.96 & 1 & 1.32 & 1687 & 0.58                   & 0.35                   & 2.38 \\ \hline
1810 & 6.52  & 1.56 & 1 & 1.43 & 1887 & 0.49                   & 0.28                   & 2.20 \\ \hline
2010 & 5.88  & 1.20 & 1 & 1.59 & 2087 & 0.40                   & 0.23                   & 1.98 \\ \hline
\end{tabular}\\
Пересчет для величин из таблицы делаем по следуюзим формулам: 
\begin{equation}
    \Theta = \frac{1}{n} \ln\left(\frac{U_i}{U_{i+n}}\right)
    \label{eq:Theta}
\end{equation}
\begin{equation}
    R_\Sigma = R + R_L
    \label{eq:Rsigma}
\end{equation}
\begin{equation}
    Q = \frac{\pi}{\Theta}
    \label{eq:Q_exper}
\end{equation}
\begin{minipage}{1\textwidth}
      \begin{tikzpicture}
        \begin{axis}[
          width=0.7\linewidth,
          height=,
          xlabel={$\frac{1}{R_\Sigma^2}\cdot 10^{-6}$, Ом$^{-2}$},
          ylabel={$\frac{1}{\Theta^2}$},
          MyStyle1,
          ]
          \draw[pink, domain=0:10] plot (\x,1.8*\x);
          \addplot [only marks] table [
                      x=1/Rsigma2, y=1/Theta2,
                      col sep=comma] {csv/2.csv};
        \end{axis}
      \end{tikzpicture}
    \end{minipage}
Из графика получаем что $k = (1.77\pm0.08) \cdot 10^6$, Ом$^2$, отсюда $R_{cr} = 2\pi \sqrt{k} = (8360\pm190)$, Ом. Полученное значение отличается от теоретического лишь на $2\frac{8360-8160}{8360+8160}\cdot 100\% \approx 2\%$.

Для следующих пунктов фиксируем два различных значения сопротивления: $R_1 = 410$ Ом, $R_2 = 2010$ Ом.

\section*{2.4 Свободные колебания на фазовой плоскости}
Фиксируем $\nu_0 = 450$ кГц и снимаем требуюмую зависимость варьируя $R$ от $R_1$ до $R_2$: \\
\begin{tabular}{|c|c|c|c|c|c|c|}
\hline
$R$, Ом    & $U_i$, В     & $U_{i+n}$, В    & $n$ & $\Theta$ & $Q_{\text{эксп}}$ & $Q_{\text{теор}}$    \\ \hline
410  & 2.8 & 1.4 & 2 & 0.35 & 9.06 & 9.96                   \\ \hline
610  & 2.5 & 0.9 & 2 & 0.51 & 6.15 & 6.69                   \\ \hline
810  & 2.2 & 1.1 & 1 & 0.69 & 4.53 & 5.04                   \\ \hline
1010 & 2.2 & 1.0 & 1 & 0.79 & 3.98 & 4.04                   \\ \hline
1210 & 4.5 & 0.6 & 2 & 1.01 & 3.12 & 3.37                   \\ \hline
1410 & 4.0 & 1.3 & 1 & 1.12 & 2.80 & 2.90                   \\ \hline
1610 & 3.6 & 1.0 & 1 & 1.28 & 2.45 & 2.54                   \\ \hline
1810 & 3.2 & 0.8 & 1 & 1.39 & 2.27 & 2.26                   \\ \hline
2010 & 2.9 & 0.6 & 1 & 1.58 & 1.99 & 2.03                   \\ \hline
\end{tabular}\\
Расчеты велись по формулам $\Theta$ - (\ref{eq:Theta}), $Q_{\text{эксп}}$(\ref{eq:Q_exper}),
\begin{equation}
    Q_{\text{теор}} = \frac{1}{R} \sqrt{\frac{C}{L}},
    \label{eq:Q_teor}
\end{equation}
где $L = 100\;\;\text{Гн}, C = 6\;\;\text{нФ}$.

\section*{2.5 Исследование резонансных кривых}
1) Фиксируем  $C = C^*$, $R = R_1$ и собираем новую схему согласно описанию. Далее ищем $\Delta \nu$ влево и вправо: 
\begin{multline*}
    U_{1\;Res} = 100\;\;\text{В} \\
    \Delta \nu_{\text{л}} = 6.45 - 5.71 = 0.74\;\;\text{кГц} \\
    \Delta \nu_{\text{п}} = 7.55 - 6.45 = 1.10\;\;\text{кГц}\\
    \delta \nu_{\text{л}} = 0.07\;\;\text{кГц}\\
    \delta \nu_{\text{п}} = 0.10\;\;\text{кГц}\\
\end{multline*}
Влево: \hspace{46.5 mm} Вправо:\\
\begin{tabular}{|c|c|c|c|}
\hline
$\nu$, Гц   & $U$, В    & $\Delta x$, мкс & $\varphi$    \\ \hline
6380 & 98.0 & 40.0 & 1.603 \\ \hline
6310 & 94.0 & 44.8 & 1.776 \\ \hline
6240 & 87.0 & 48.4 & 1.898 \\ \hline
6170 & 79.2 & 53.2 & 2.062 \\ \hline
6100 & 71.2 & 57.0 & 2.185 \\ \hline
6030 & 63.6 & 60.0 & 2.273 \\ \hline
5960 & 57.6 & 63.4 & 2.374 \\ \hline
5890 & 51.6 & 65.6 & 2.428 \\ \hline
5820 & 45.4 & 67.8 & 2.479 \\ \hline
5750 & 41.6 & 69.8 & 2.522 \\ \hline
\end{tabular}
\begin{tabular}{|c|c|c|c|}
\hline
$\nu$, Гц   & $U$, В    & $\Delta x$, мкс & $\varphi$    \\ \hline
6550 & 98.8 & 28.0 & 1.152 \\ \hline
6650 & 91.6 & 22.2 & 0.928 \\ \hline
6750 & 82.8 & 18.1 & 0.768 \\ \hline
6850 & 74.4 & 14.6 & 0.628 \\ \hline
6950 & 66.4 & 12.2 & 0.533 \\ \hline
7050 & 60.4 & 10.3 & 0.456 \\ \hline
7150 & 55.2 & 8.9  & 0.400 \\ \hline
7250 & 50.8 & 7.6  & 0.346 \\ \hline
7350 & 46.8 & 6.7  & 0.309 \\ \hline
7450 & 43.2 & 5.8  & 0.271 \\ \hline
7550 & 41.2 & 5.4  & 0.256 \\ \hline
\end{tabular}\\
2) Теперь меняем сопротивление магазина на $R_2$
\begin{multline*}
    U_{2\;Res} = 26.4\;\;\text{В} \\
    \nu_0 = 6.9\;\;\text{кГц} \\
    \Delta \nu_{\text{л}} = 6.9 - 4.4 = 2.5\;\;\text{кГц} \\
    \Delta \nu_{\text{п}} = 12.2 - 6.9 = 5.3\;\;\text{кГц}\\
    \delta \nu_{\text{л}} = 250\;\;\text{Гц}\\
    \delta \nu_{\text{п}} \neq const\\
\end{multline*}
Влево: \hspace{46.5 mm} Вправо:\\
\begin{tabular}{|c|c|c|c|}
\hline
$\nu$, Гц   & $U$, В    & $\Delta x$, мкс & $\varphi$    \\ \hline
6900 & 26.6 & 19.3 & 0.837 \\ \hline
6650 & 26.6 & 22.7 & 0.948 \\ \hline
6400 & 26.0 & 27.0 & 1.086 \\ \hline
6150 & 24.8 & 32.2 & 1.244 \\ \hline
5900 & 22.6 & 36.9 & 1.368 \\ \hline
5650 & 20.4 & 42.1 & 1.495 \\ \hline
5400 & 18.4 & 47.5 & 1.612 \\ \hline
5150 & 16.4 & 52.4 & 1.696 \\ \hline
4900 & 14.4 & 57.8 & 1.780 \\ \hline
4650 & 12.3 & 62.8 & 1.835 \\ \hline
4400 & 10.6 & 68.3 & 1.888 \\ \hline
\end{tabular}
\begin{tabular}{|c|c|c|c|}
\hline
$\nu$, Гц   & $U$, В    & $\Delta x$, мкс & $\varphi$ \\ \hline
7200  & 26.2 & 15.5 & 0.701 \\ \hline
7500  & 25.8 & 12.5 & 0.589 \\ \hline
7800  & 24.8 & 10.0 & 0.490 \\ \hline
8100  & 23.0 & 8.3  & 0.422 \\ \hline
8400  & 22.4 & 6.7  & 0.354 \\ \hline
8700  & 21.6 & 5.5  & 0.301 \\ \hline
9000  & 20.8 & 4.6  & 0.260 \\ \hline
9300  & 20.0 & 3.8  & 0.222 \\ \hline
9600  & 19.2 & 3.1  & 0.187 \\ \hline
10600 & 17.2 & 1.6  & 0.109 \\ \hline
11600 & 16.3 & 0.9  & 0.064 \\ \hline
12600 & 15.5 & 0.4  & 0.035 \\ \hline
\end{tabular}\\
\begin{minipage}{1\textwidth}
      \begin{tikzpicture}
        \begin{axis}[
          width=0.8\linewidth,
          height=,
          xlabel={$\frac{\nu}{\nu_0}$},
          ylabel={$\frac{U}{U_0}$},
          MyStyle1,
          ]
          \draw[blue, domain=0:10] plot (\x,0.707106);
          \addplot [only marks] table [
                      x=1nu/nu0, y=1U/U0,
                col sep=comma] {csv/7.csv};
          \addplot [purple, only marks] table [
                      x=2nu/nu0, y=2U/U0,
                col sep=comma] {csv/7.csv};
        \end{axis}
      \end{tikzpicture}
    \end{minipage}
Из данного графика резонансной кривой получаем
\begin{equation*}
    \begin{cases}
        Q_{R_1} = \frac{\omega_0}{2\Delta\Omega} = \frac{\nu_0}{2 \Delta \nu} = 9.1 \\
        \Delta Q_{R_1} = Q_{R_1} \cdot \frac{\delta(\Delta \nu)}{\Delta \nu} = 0.2
    \end{cases} \Rightarrow
    Q_{R_1} = (9.12\pm0.07)
\end{equation*}
\begin{equation*}
    \begin{cases}
        Q_{R_2} = \frac{\omega_0}{2\Delta\Omega} = \frac{\nu_0}{2 \Delta \nu} = 1.5 \\
        \Delta Q_{R_2} = Q_{R_2} \cdot \frac{\delta(\Delta \nu)}{\Delta \nu} = 0.02
    \end{cases} \Rightarrow
    Q_{R_2} = (1.50\pm0.02)
\end{equation*}
Теперь строим ФЧХ по полученным данным:\\
\begin{minipage}{1\textwidth}
      \begin{tikzpicture}
        \begin{axis}[
          width=0.8\linewidth,
          height=,
          xlabel={$\nu$, кГц},
          ylabel={$\varphi$},
          MyStyle1,
          ]
          \draw[blue, domain=0:100] plot (\x,1.5707963);
          \draw[green, domain=0:100] plot (\x,2.35619);
          \addplot [only marks] table [
                      x=1nu, y=1phi,
                 col sep=comma] {csv/9.csv};
          \addplot [only marks] table [
                      x=1nu, y=sim1phi,
                 col sep=comma] {csv/9.csv};
          \addplot[purple, only marks] table [
                      x=2nu, y=2phi,
                 col sep=comma] {csv/9.csv};
          \addplot[purple, only marks] table [
                      x=2nu, y=sim2phi,
                 col sep=comma] {csv/9.csv};    
        \end{axis}
      \end{tikzpicture}
    \end{minipage}
Из графика видно чток сожалению для $R_2$ у нас не выйдет определить добротность, а для $R_1$:
\begin{equation*}
    \begin{cases}
        Q_{R_1} = \frac{\omega_0}{\Delta \omega} = \frac{6.45} {0.75} = 8.6 \\
        \Delta Q_{R_1} = Q_{R_1} \cdot \frac{\delta(\Delta \nu)}{\Delta \nu} = 0.2
    \end{cases} \Rightarrow
    Q_{R_1} = (8.6\pm0.2)
\end{equation*}

\section*{2.6 Процессы установления и затухания}
Выставляем $R = R_1$, генератор в режим Burst $\nu_0 = 6.5$ кГц, $N = 15$. 
\begin{figure}[h]
\centering
\includegraphics[width=0.5\linewidth]{3.2.4 Свободные и вынужденные колебания в электрическом контуре/картинки/цуги.png}
\caption{Примерный рисунок на осциллографе}
\label{fig:czugi}
\end{figure}
По картине на осциллографе получаем следующие данные:
\begin{multline*}
    U_0 = 99.6\;\;\text{В}\\
    \textbf{1}:\;\; U_k = 33.6\;\;\text{В},\;U_{k+4} = 84.4\;\;\text{В} \Rightarrow \Theta = 0.36 \\
    \textbf{2}:\;\; U_k = 54.4\;\;\text{В},\;U_{k+6} = 94.8\;\;\text{В} \Rightarrow \Theta = 0.34 \\
    \textbf{3}:\;\; U_m = 70.4\;\;\text{В},\;U_{m+2} = 34.4\;\;\text{В} \Rightarrow \Theta = 0.35 \\
\end{multline*}
Для сопротивления $R = R_2$, получаем:
\begin{multline*}
    U_0 = 26.8\;\;\text{В}\\
    \textbf{1}:\;\; U_k = 9.0\;\;\text{В},\;U_{k+2} = 26.4\;\;\text{В} \Rightarrow \Theta = 1.90\\
    \textbf{2}:\;\; U_m = 26.6\;\;\text{В},\;U_{m+1} = 21.0\;\;\text{В} \Rightarrow \Theta = 3.36\\
\end{multline*}
Где $\Theta = \frac{1}{n} \ln \left( \frac{U_0-U_k}{U_0-U_{k+n}}\right)$.

Определим добротность при $R_1$ и $R_2$, нарастании и затухания как: $Q_i = \frac{\pi}{<\Theta_i>}$, а $\Delta Q_1 = Q_1 \cdot \frac{\sigma(\Theta)}{<\Theta>} \rightarrow$
\begin{equation*}
    Q_1^{\text{нарастание}} = (9.0\pm0.3)
\end{equation*}
\begin{equation*}
    Q_1^{\text{затухание}} = (7.9\pm0.8)
\end{equation*}
\begin{equation*}
    Q_2^{\text{нарастание}} = (1.65\pm0.6)
\end{equation*}
\begin{equation*}
    Q_2^{\text{затухание}} = (13.6\pm4)
\end{equation*}

\section*{2.7 Обработка результатов}
Почти вся в прошлых пунктах. \\
\begin{tabular}{|c|ccc|cccc|}
\hline
\multirow{2}{*}{$R$} & \multicolumn{3}{c|}{Свободные колебания}                  & \multicolumn{4}{c|}{Вынужденные колебания}                                                        \\ \cline{2-8} 
                   & \multicolumn{1}{c|}{$f(L,C,R)$} & \multicolumn{1}{c|}{$f(\Theta)$} & Спираль & \multicolumn{1}{c|}{АЧХ} & \multicolumn{1}{c|}{ФЧХ} & \multicolumn{1}{c|}{Нарастание} & Затухание \\ \hline
$R_1$ = & \multicolumn{1}{c|}{10.0}  & \multicolumn{1}{c|}{8.5}  &    9.1     & \multicolumn{1}{c|}{$(9.12\pm0.07)$}    & \multicolumn{1}{c|}{(8.6\pm0.2)}    & \multicolumn{1}{c|}{$(9.0\pm0.3)$}           &     $(7.9\pm0.8)$      \\ \hline
$R_2$ = & \multicolumn{1}{c|}{2.0}  & \multicolumn{1}{c|}{2.0}  &     2.0    & \multicolumn{1}{c|}{$(1.50\pm0.02)$}    & \multicolumn{1}{c|}{-}    & \multicolumn{1}{c|}{$(1.65\pm0.6)$}           &      -     \\ \hline
\end{tabular}





\end{document}